\itemhead{复习题一 · 第3题（教材p.29）}{XJ-B1-C01-RV-P029-Y0779}
\label{XJ-B1-C01-RV-P029-Y0779}
{\footnotesize 湘教版教材p.29 / PDF第36页 / 复习题一 / 第3题}
\field{原题}{（多选题）已知集合 $N=\{m+3,m^2+2m,(m-1)^2\}$，若 $-1$ 是集合 $N$ 的元素，则 $m$ 的值是（\quad）
\par A. $-4$\qquad B. $-1$\qquad C. $-2$\qquad D. $0$}
\field{参考答案}{AB。由 $m+3=-1$ 得 $m=-4$；由 $m^2+2m=-1$ 得 $m=-1$；平方不可能为 $-1$，且两候选均符合元素互异性。}
{\small \field{关联知识}{S01 集合、元素与常用数集（教材帮p.1、2 / PDF10、11）}
\field{方法与易错点}{先辨认研究对象和数集范围，再代入定义判断；含参集合逐个来源分类。 易错：把自然数集误作正整数集；混淆元素和单元素集合；候选参数造成元素重复。}
\field{答案依据}{答案册相关题组：PDF第8、9页（册内页7、8）；按当前题库原题逐项重算，不按旧版题号机械对套。}}
